\section{April 2}


\subsection{Compatible almost complex structures}

  \begin{proof}[Proof of \cref{thm:compatible_almost_complex_structures_and_riemannian_metrics} (continued)]
    Given a Riemannian metric $g$ on $M$, since both $g$ and $\omega$ define non-degenerate bilinear forms on each tangent space, there exists a unique bundle map $A: TM \to TM$ such that $\omega(X, Y) = g(AX, Y)$ for all vector fields $X, Y$ on $M$.

    Fix $p \in M$ and consider $A|_p: T_p M \to T_p M$ for each $p \in M$.
    Let $A^*|_p$ be the adjoint of $A|_p$ with respect to $g$, i.e., $g(A^*|_p X, Y) = g(X, A|_p Y)$ for all $X, Y \in T_p M$.
    We have $A^*|_p = -A|_p$ since 
    \begin{align*}
      g(A^*|_p X, Y) & = g(X, A|_p Y) \\
        & = g(AY, X) = \omega(Y, X) \\
        & = -\omega(X, Y) = -g(AX, Y).
    \end{align*}
    Let $Q = A A^* = -A^2$, then $Q$ is a positive definite self-adjoint operator on $T_p M$.
    Then $Q$ has a unique positive definite self-adjoint square root $P$ such that $P^2 = Q$.

    \begin{claim}\label{claim:AP_PA_in_the_proof_of_compatible_almost_complex_structures_and_riemannian_metrics}
      $AP = PA$.
    \end{claim}

    Define $\Phi(g) := J|_p = A|_pP^{-1} $.
    Then $J|_p^2 = A|_p P^{-1} A|_p P^{-1} = A|_p A|_p P^{-2} = -\id$.
    Let $g_J(u,v) := \omega(u, J|_p v)$ for $u, v \in T_p M$.
    We have 
    \begin{align*}
      g_J(u, v) & = \omega(u, J|_p v) = g(u, A|_p P^{-1} v) = g(A|_p u, AP^{-1} v) \\
        & = g(u, A^* A P^{-1} v) = g(u, P^2 P^{-1} v) = g(u, Pv).
    \end{align*}
    Since $P$ is self-adjoint and positive definite, $g_J$ is a Riemannian metric on $T_p M$.
    Hence $J|_p$ is compatible with $\omega$.
    In particular, we get a map $\Phi: \calG \to \calJ_\omega$.

    Consider the composition $\calJ_\omega \to \calG \xrightarrow{\Phi} \calJ_\omega$.
    Given $J \in \calJ_\omega$, define $g_J(X, Y) = \omega(X, JY)$ for all $X, Y \in T_p M$.
    Note that $g_J(J(u),v) = g_J(v, J(u)) = \omega(v, J^2 u) = \omega(u, v)$ for all $u, v \in T_p M$.
    Hence we can take $A = J$ and $P = \id$ in the construction of $\Phi(g_J)$, which implies $\Phi(g_J) = J$.
  \end{proof}

  \begin{proof}[Proof of \cref{claim:AP_PA_in_the_proof_of_compatible_almost_complex_structures_and_riemannian_metrics}]
    Since $A^* = -A$, there exists an orthonormal basis such that $A$ has the block diagonal form
    \[ A = \begin{pmatrix}
        0 & a_1 & & & \\
        -a_1 & 0 & & & \\
        & & \ddots & & \\
        & & & 0 & a_n \\
        & & & -a_n & 0
      \end{pmatrix} \]
    where $a_i > 0$ for $i = 1, \dots, n$.
    Then $Q = A A^* = -A^2$ has the block diagonal form
    \[ Q = \begin{pmatrix}
        a_1^2 & & & & \\
        & a_1^2 & & & \\
        & & \ddots & & \\
        & & & a_n^2 & \\
        & & & & a_n^2
      \end{pmatrix} \]
    and $P = Q^{1/2}$ has the block diagonal form
    \[ P = \begin{pmatrix}
        a_1 & & & & \\
        & a_1 & & & \\
        & & \ddots & & \\
        & & & a_n & \\
        & & & & a_n
      \end{pmatrix}. \]
    It is easy to see that $AP = PA$.
  \end{proof}

  \begin{remark}
    Endow $\calJ_\omega$ with suitable topology.
    Then it is a $\infty$-dimensional Fréchet manifold. 
  \end{remark}

  \begin{theorem}
    The space $\calJ_\omega$ of all compatible almost complex structures on $(M, \omega)$ is contractible.
  \end{theorem}
  \begin{proof}
    Denote by $\Psi: \calJ_\omega \to \calG$ the map defined by $J \mapsto g_J$ where $g_J(X, Y) = \omega(X, JY)$ for all $X, Y \in T_p M$.
    Consider the map
    \[ F: [0,1] \times \calJ_\omega \to \calJ_\omega, \quad F(t, J) = \Phi((1-t) \Psi(J) + t \Psi(J_0)) \]
    where $J_0$ is a fixed compatible almost complex structure on $(M, \omega)$.
    Then $F(0, J) = \Phi(\Psi(J)) = J$ and $F(1, J) = \Phi(\Psi(J_0)) = J_0$ for all $J \in \calJ_\omega$.
  \end{proof}

  \begin{corollary}
    Any two compatible almost complex structures on $(M, \omega)$ can be connected by a smooth path in $\calJ_\omega$.
  \end{corollary}


\subsection{Symplectic group actions}
  
  \begin{definition}
    Let $G$ be a Lie group acting on a symplectic manifold $(M, \omega)$.
    We say the action is \textit{symplectic} if $g^* \omega = \omega$ for all $g \in G$.
  \end{definition}

  Fix a Lie group $G$ acting symplectically on $(M, \omega)$.
  Let $\frakg$ be the Lie algebra of $G$.
  For $X \in \frakg$, it induces a vector field $\tilde{X}$ on $M$ defined by
  \[ \tilde{X}_p = \frac{d}{dt}\Big|_{t=0} \exp(tX) \cdot p, \quad p \in M. \]
  One can show that $\tilde{X}$ is a symplectic vector field, i.e., $i_{\tilde{X}} \omega $ is closed.

  Recall that we have 
  \[ [\tilde{X}, \tilde{Y}] = -\widetilde{[X, Y]} \]
  for all $X, Y \in \frakg$.
  In other words, the map $\frakg \to SV(M, \omega), X \mapsto \tilde{X}$ is a Lie algebra anti-homomorphism.

  Let $SV(M, \omega)$ (resp. $HV(M, \omega)$) be the space of all symplectic vector fields (resp. Hamiltonian vector fields) on $(M, \omega)$.
  We have 
  \begin{equation}\label{eq:commutative_diagram_for_symplectic_group_actions}
    \begin{tikzcd}
      \frakg \ar[d] \ar[r, gray, "\tilde{\mu}"] & C^\infty(M) \ar[d] \\
      SV(M, \omega) & HV(M, \omega). \ar[l, "\supseteq"']
    \end{tikzcd}
  \end{equation}
  We will ask: is there a Lie algebra anti-homomorphism $\tilde{\mu}: \frakg \to C^\infty(M)$ such that the above diagram commutes?

  \begin{definition}
    If such a Lie algebra anti-homomorphism $\tilde{\mu}: \frakg \to C^\infty(M)$ exists, we call the action of $G$ on $(M, \omega)$ \textit{Hamiltonian} and call $\tilde{\mu}$ a \textit{comoment map}.

    Given a comoment map $\tilde{\mu}: \frakg \to C^\infty(M)$, we can define a moment map $\mu: M \to \frakg^*$ by $\langle \mu(p), X \rangle = \tilde{\mu}(X)(p)$ for all $p \in M$ and $X \in \frakg$.
    % Then $\mu$ is $G$-equivariant with respect to the given action of $G$ on $M$ and the coadjoint action of $G$ on $\frakg^*$.
  \end{definition}

  \begin{lemma}
    Assume $G$ is commutative and $M$ is compact.
    If there exists a linear map $\tilde{\mu}: \frakg \to C^\infty(M)$ such that \cref{eq:commutative_diagram_for_symplectic_group_actions} commutes, then $\tilde{\mu}$ is a comoment map.
  \end{lemma}
  \begin{proof}
    Since $G$ is commutative, we have $[X, Y] = 0$ for all $X, Y \in \frakg$.
    Hence $\tilde{\mu}([X, Y]) = 0$ for all $X, Y \in \frakg$.

    We have 
    \[ \text{Hamiltonian vector field of } {\{\tilde{\mu}(X), \tilde{\mu}(Y)\}} = [\tilde{X}, \tilde{Y}] = -\widetilde{[X, Y]} = 0. \]
    Then $d \{\tilde{\mu}(X), \tilde{\mu}(Y)\} = 0$ for all $X, Y \in \frakg$.
    Hence $\{\tilde{\mu}(X), \tilde{\mu}(Y)\}$ is a constant function on $M$ for all $X, Y \in \frakg$.
    Note that $\{\tilde{\mu}(X), \tilde{\mu}(Y)\} = \tilde{X}(\tilde{\mu}(Y)) = 0$ at any critical point of $\tilde{\mu}(Y)$.
    Since $M$ is compact, $\tilde{\mu}(Y)$ has a critical point.
    Hence $\{\tilde{\mu}(X), \tilde{\mu}(Y)\} = 0$ for all $X, Y \in \frakg$.
  \end{proof}

  \begin{lemma}
    Let $G$ acts on $(M, \omega)$ Hamiltonianly with a moment map $\mu: M \to \frakg^*$.
    Consider the Poisson structure on $\frakg^*$ and $M$.
    Then $\mu$ is a anti-Poisson map, i.e., 
    \[ \{f \circ \mu, g \circ \mu\} = -\{f, g\} \circ \mu \]
    for all $f, g \in C^\infty(\frakg^*)$.
  \end{lemma}
  \begin{proof}
    Exercise.
  \end{proof}